\section{Discussion}
\label{sec:discussion}

The main value of the workflow is that it makes several failure modes explicit. A low extremal-index estimate may indicate clustered extremes under the chosen threshold and run length, but it can also result from periodicity, persistent operating regimes, preprocessing, missingness, or finite-sample noise. Dynamical EVT supplies the language linking recurrence and extremes, yet the industrial interpretation still requires domain evidence. This paper therefore treats \(\theta\) as a dependence diagnostic and not as a direct mechanical degradation variable.

Event-level evaluation changes the operational reading of anomaly scores. Pointwise score improvements can be unhelpful if they fragment one failure into many duplicate alarms or if they achieve recall only by keeping a machine under warning for a large fraction of exposure. Conversely, an apparently sparse alarm policy can be unacceptable if it fires too late for maintenance action. The generated frontier and timeline make these trade-offs visible, but they do not remove the need for cost, staffing, and repair-window assumptions.

The five-dataset matrix changes the evidential status of the package. The repository is no longer only a synthetic workflow plus one compressor-family case study. It now verifies public raw and processed artifacts for compressor telemetry, truck repair histories, hydraulic test-rig cycles, and semiconductor process examples, and it writes one terminal diagnostic row for each. This broader base is valuable because it forces the workflow to handle different schemas, labels, missingness patterns, entity structures, and independence units.

The negative result in CLM-007 is still important. The current generated assets do not support broad industrial performance claims beyond the evaluated evidence. That statement is not a weakness to hide; it is a boundary condition for honest use of the artifacts. MetroPT and MetroPT2 have limited independent failure evidence; SCANIA changes the target estimand to vehicle-level risk stratification; Hydraulic Systems uses laboratory cycle-level component states; and SECOM uses wafer-level yield labels with missing high-dimensional process measurements. Future empirical work should therefore add additional public or prospective datasets and predeclared baselines before making a broader predictive-maintenance claim.

The reproducibility layer is also a scientific constraint. Every generated figure and table carries provenance, the claim ledger blocks unsupported manuscript injection, and the paper-check command audits citations, claims, generated assets, and compiled PDFs. This structure does not guarantee that every modelling choice is optimal, but it makes unsupported interpretation harder to introduce silently.
